% Bewertungspaket zum Gesamttest «Zentrische Streckung und Ähnlichkeit» — Quelle des PDFs
% (python3 scripts/build-lp-pdf.py aehnlichkeit). Ganze Punkte je Teilschritt; Werte mit python3 nachgerechnet
% (08.10.2026, scripts/lp/aehnlichkeit/zahlen.py, Folgefehler-Fälle eingeschlossen). Aufbau wie trigonometrische-berechnungen/bewertungspaket.tex.
\documentclass[11pt]{article}
\usepackage{../lp-druck}
\renewcommand{\lpname}{Zentrische Streckung und Ähnlichkeit}
\renewcommand{\lpbereich}{Grundlagenfach}
\renewcommand{\lpteilgebiet}{5.2}
\renewcommand{\lpfuss}{Bewertungspaket}
\newcolumntype{P}{>{\centering\arraybackslash}p{10mm}}
\newenvironment{raster}{\par\smallskip\noindent\tabularx{\linewidth}{@{}>{\raggedright\arraybackslash}p{38mm}P X@{}}\toprule
  \small\textsc{Teilschritt} & \small\textsc{P} & \small\textsc{Musterlösung}\\\midrule}{\bottomrule\endtabularx\par}
\newcommand{\fehler}[1]{\par\smallskip{\small\textcolor{lprot}{\bfseries Typische Fehler:} #1}}
\newcommand{\E}{\,{\footnotesize\bfseries\color{lpgruen}(E)}}
\newcommand{\A}{\,{\footnotesize\bfseries\color{lpblau}(A)}}
\begin{document}
\lpkopf{Bewertungspaket zum Gesamttest}{}

\begin{lpachtung}{Erst lösen, dann öffnen}
Dieses Paket enthält die Musterlösung. Wer es vor dem Lösen liest, misst am Ende nur, wie gut das Abschreiben gelingt.
\end{lpachtung}

\section*{1 · So gehst du vor}
\begin{enumerate}[leftmargin=*,itemsep=0pt]
\item \textbf{Gesamttest lösen} — vollständig, mit Skizze und Rechenweg, ohne dieses Paket; der Taschenrechner ist erlaubt.
\item \textbf{Fotografieren oder scannen}, gut lesbar, eine Seite pro Bild. \textbf{Kein Name} auf den Bildern.
  Handyfotos tragen oft unsichtbar Standort, Zeit und Gerät mit (Metadaten). Lade darum lieber einen Scan oder ein
  Bildschirmfoto des Fotos hoch, oder entferne vorher die Standortangaben.
\item \textbf{Bewerten lassen.} Ob du dafür eine KI nutzt und welche, entscheidest du selbst und in eigener Verantwortung —
  je nachdem, was dir aufgrund deines Alters und deiner persönlichen Situation erlaubt ist (Altersgrenzen und
  Nutzungsbedingungen des Anbieters, allenfalls Einverständnis deiner Eltern). Lade nur deine Lösung und dieses PDF hoch,
  keine persönlichen Daten, und füge den Auftrag aus Abschnitt~2 ein. Ohne KI kannst du dich mit Abschnitt~3 selbst bewerten.
\item \textbf{Rückmeldung prüfen.} Stimmt, was die KI aus deiner Handschrift gelesen hat? Eine KI kann sich irren —
  im Zweifel gilt das Raster in Abschnitt~3.
\item \textbf{Gezielt wiederholen} nach Abschnitt~4.
\end{enumerate}

\needspace{8\baselineskip}
\section*{2 · Auftrag an die KI}
\begin{tcolorbox}[breakable,colback=lplinie!25,colframe=lplinie,boxrule=0.5pt,arc=1mm,fontupper=\small,left=2mm,right=2mm,top=1.5mm,bottom=1.5mm]
Du bist Korrektorin oder Korrektor für Mathematik an einer Schweizer Berufsmaturitätsschule. Im Anhang findest du (1)~das Bewertungspaket zum Gesamttest «Zentrische Streckung und Ähnlichkeit» mit Musterlösung und Punkteraster und (2)~Fotos meiner handschriftlichen Lösung.\par\smallskip
Geh so vor:\par
1. Schreib zu jeder Aufgabe G1 bis G7 zuerst ab, was du in meiner Lösung liest — Rechenweg und Ergebnis. Was du nicht sicher lesen kannst, markierst du mit [unsicher]. Nicht raten.\par
2. Vergib die Punkte genau nach dem Raster im Bewertungspaket (Abschnitt 3), ganze Punkte je Zeile. Ziehe einen Fehler nur einmal ab: Rechne ich mit einem falschen Zwischenergebnis richtig weiter, gibt es die Folgepunkte — aber nur, wenn der Schritt dieser Zeile selbst richtig ist. Mache ich denselben Fehler in einer weiteren Zeile nochmals (zum Beispiel wieder eine Fläche mit dem Längenfaktor), gibt auch diese Zeile 0. Kein Folgepunkt für ein Ergebnis, das in der Sache unmöglich ist (eine negative Länge, ein Bild, das grösser ist als die Schachtel). Zeilen mit \textbf{(E)} gibt es nur für das richtige Ergebnis, nie als Folgepunkt. Die Zeilen «Typische Fehler» sagen, welche Rasterzeile bei diesem Fehler 0 Punkte gibt — sie sind kein zusätzlicher Abzug.\par
3. Andere richtige Lösungswege zählen voll, ebenso gleichwertige Schreibweisen: Dezimalpunkt oder Dezimalkomma, exakte Werte wie $1.6\sqrt{3}$ statt Dezimalzahlen, anders, aber vernünftig gerundete Ergebnisse ($\pm 0.01$ oder durch gerundete Zwischenwerte erklärbar). \textbf{Einheiten:} Fehlt bei einem Endergebnis die Einheit oder ist sie falsch ($\text{cm}$ statt $\text{cm}^2$), gibt die (E)-Zeile dieses Ergebnisses 0 Punkte — aber nur beim ersten Mal im ganzen Gesamttest; weitere Einheitenfehler kosten nichts mehr. Verlangt die Aufgabe eine Einheit («in km»), gehört die Umrechnung zum Ergebnis. Gemischte Einheiten in einer Rechnung (Meter und Zentimeter in einer Verhältnisgleichung) sind kein Einheitenfehler, sondern ein Rechenfehler: Die Zeile gibt 0. Zeilen mit \textbf{(A)} verlangen nur Erkennen oder Ablesen, diese Punkte gibt es auch ohne Rechenweg. Zeilen mit \textbf{(E)} sind Ergebnispunkte: Das Ergebnis muss stimmen \emph{und} aus einem erkennbaren Weg hervorgehen. Alle übrigen Zeilen bewerten den Weg selbst.\par
4. Begründe jeden Abzug in einem Satz und nenne die Fehlerart (zum Beispiel «Fläche mit $k$ statt mit $k^2$ gerechnet» oder «$AA'$ statt $SA'$ im 2. Strahlensatz»). Schreib die Musterlösung nicht ab — zeig mir, wo mein Fehler liegt.\par
5. Gib zum Schluss eine Tabelle «Aufgabe | Punkte | Begründung», die Gesamtpunktzahl (von 25) und — nach der Selbsteinschätzung in Abschnitt 4 — welches Kapitel des Leitprogramms ich wiederholen soll. Keine Note.\par
6. Fehlt ein Foto oder ist eine Aufgabe unlesbar, sag es, statt sie zu bewerten.
\end{tcolorbox}

\needspace{8\baselineskip}
\section*{3 · Musterlösung und Punkteraster}
\textbf{Allgemein:} Ganze Punkte je Zeile. Ein Fehler wird nur einmal abgezogen (Folgefehler: wer mit einem falschen
Zwischenergebnis richtig weiterrechnet, bekommt die Wegpunkte der folgenden Zeilen, aber keine (E)-Punkte). \textbf{Eine Regel
für Folgepunkte:} Es gibt sie nur, wenn der Schritt der Zeile selbst richtig ist und nur ein Wert aus einer früheren Zeile falsch
übernommen wird. Wer denselben Fehler nochmals macht (wieder Fläche mit dem Längenfaktor, wieder nach Buchstaben zugeordnet),
bekommt auch in dieser Zeile 0. Kein Folgepunkt für ein in der Sache unmögliches Ergebnis (negative Länge; ein Bild, das grösser
ist als die Schachtel). In der Spalte P:
\textbf{(A)}~= erkennen oder ablesen, zählt auch ohne Rechenweg · \textbf{(E)}~= Ergebnispunkt, Ergebnis richtig
\emph{und} Weg erkennbar · ohne Zeichen = die Zeile bewertet den Weg. Taschenrechner erlaubt; gerundete Ergebnisse auf zwei
Dezimalen, Toleranz $\pm 0.01$ (oder durch gerundete Zwischenwerte erklärbar). \textbf{Einheiten:} Der erste Einheitenfehler im
ganzen Test kostet die (E)-Zeile dieses Ergebnisses, weitere nichts; verlangt die Aufgabe eine Einheit, gehört die Umrechnung zum
Ergebnis; gemischte Einheiten in einer Rechnung sind ein Rechenfehler (Zeile 0). \textbf{Längen sind positiv:} Eine negative Länge
(aus $k < 0$) ist ein Fehler, nicht eine gleichwertige Schreibweise.
«Typische Fehler» nennt, welche Zeile 0 Punkte gibt, und die Punkte, die bleiben.
\textbf{Teilrichtige Zeilen:} Verlangt eine Zeile mehrere Angaben und stimmt nur ein Teil, gibt die Zeile 0 Punkte.

\begin{lpaufgabe}{G1}{Zentrische Streckung mit Koordinaten}{3 P}
\begin{raster}
(a) $P'$ & 1\E & Weg $Z \to P$: $(2 \mid 1)$; mal $-1.5$: $(-3 \mid -1.5)$; $P' = (2 - 3 \mid -1 - 1.5) = (-1 \mid -2.5)$\\
(b) $\overline{P'Q'}$ & 1\E & $|k| \cdot \overline{PQ} = 1.5 \cdot \sqrt{13} \approx 5.41\,\text{cm}$ (mit $3.61$ gerechnet: $5.42$, zählt)\\
(c) parallel & 1 & $Q'R'$, das Bild von $QR$: Bei einer zentrischen Streckung ist jede Bildseite parallel zu ihrer Originalseite (auch bei $k < 0$). Ohne Begründung: 0.\\
\end{raster}
\fehler{Koordinaten von $P$ mit $-1.5$ multipliziert, $P'(-6 \mid 0)$ (vom Ursprung aus gestreckt): (a) 0; (b) und (c) zählen, wenn richtig $\to$ 2 von 3\,P.
Vorzeichen von $k$ vergessen, $P'(5 \mid 0.5)$ (Bild auf derselben Seite von $Z$): (a) 0 $\to$ höchstens 2 von 3\,P.
$\overline{P'Q'} = -5.41$ (negative Länge) oder $2.25 \cdot \sqrt{13} \approx 8.11$ (mit $k^2$): (b) 0.}
\end{lpaufgabe}

\begin{lpaufgabe}{G2}{Strahlensätze}{4 P}
\begin{raster}
(a) Ansatz & 1 & 1. Strahlensatz, z.\,B. $\overline{SA} : \overline{AA'} = \overline{SB} : \overline{BB'}$, also $\tfrac{2.5}{3.5} = \tfrac{3}{\overline{BB'}}$ (gleichwertig über $\overline{SA'} = 6$ und $\overline{SB'} = 7.2$)\\
(a) $\overline{BB'}$ & 1\E & $\overline{BB'} = \tfrac{3 \cdot 3.5}{2.5} = 4.2\,\text{cm}$\\
(b) Ansatz & 1 & 2. Strahlensatz mit den ganzen Strecken ab $S$: $\overline{SA'} = 2.5 + 3.5 = 6$, $\tfrac{\overline{A'B'}}{\overline{AB}} = \tfrac{\overline{SA'}}{\overline{SA}}$\\
(b) $\overline{A'B'}$ & 1\E & $\overline{A'B'} = 2 \cdot \tfrac{6}{2.5} = 4.8\,\text{cm}$\\
\end{raster}
\fehler{$\overline{A'B'} = 2 \cdot \tfrac{3.5}{2.5} = 2.8$ ($\overline{AA'}$ statt $\overline{SA'}$): beide (b)-Zeilen 0 $\to$ (a) zählt, 2 von 4\,P.
$\overline{SB'} = 7.2$ als $\overline{BB'}$ angegeben: (a)-Ergebnis 0, Ansatz zählt.
Verhältnis verkehrt ($\overline{BB'} = 3 \cdot \tfrac{2.5}{3.5} \approx 2.14$; $\overline{A'B'} = 2 \cdot \tfrac{2.5}{6} \approx 0.83$): die Ansatz- und Ergebniszeile der Teilaufgabe 0.}
\end{lpaufgabe}

\begin{lpaufgabe}{G3}{Lochkamera}{4 P}
\begin{raster}
(a) Skizze, Begründung & 1 & Zwei Strahlen (Kopf, Füsse) kreuzen sich im Loch $S$: X-Figur. Mensch und Bild stehen beide senkrecht, sind also parallel — darum gelten die Strahlensätze. Die Abstände von $S$ zu den beiden Parallelen stehen im selben Verhältnis (die Streckung mit Zentrum $S$ bildet auch das Lot ab). Skizze \emph{und} Begründung (Parallelen) nötig.\\
(b) Ansatz & 1 & $\tfrac{b}{1.8\,\text{m}} = \tfrac{0.15\,\text{m}}{6\,\text{m}}$: Bild zu Mensch wie Abstand zu Abstand (gleichwertig mit $15\,\text{cm}$ und $600\,\text{cm}$)\\
(b) Bildgrösse & 1\E & $b = 1.8 \cdot \tfrac{0.15}{6} = 0.045\,\text{m} = 4.5\,\text{cm}$\\
(c) Abstand & 1\E & $\tfrac{0.03}{1.8} = \tfrac{0.15}{d}$, $d = \tfrac{1.8 \cdot 0.15}{0.03} = 9\,\text{m}$\\
\end{raster}
\fehler{Einheiten gemischt ($b = 1.8 \cdot \tfrac{15}{6} = 4.5$, als Meter; $d = \tfrac{1.8 \cdot 0.15}{3} = 0.09$): die Ergebniszeile 0 — auch wenn schon ein Einheitenfehler gezählt wurde: Ein $4.5\,\text{m}$ grosses Bild in einer Schachtel ist unmöglich. Der Ansatz zählt, wenn das Verhältnis stimmt.
Verhältnis verkehrt ($b = 1.8 \cdot \tfrac{6}{0.15} = 72\,\text{m}$): beide (b)-Zeilen 0; (c) zählt, wenn richtig $\to$ höchstens 2 von 4\,P.
(a) nur Skizze, ohne Parallelen: (a) 0.}
\end{lpaufgabe}

\begin{lpaufgabe}{G4}{Karte}{4 P}
\begin{raster}
(a) Weg & 1\E & $18.4\,\text{cm} \cdot 25\,000 = 460\,000\,\text{cm} = 4.6\,\text{km}$\\
(b) Faktor & 1 & Flächen wachsen mit $25\,000^2 = 625\,000\,000$\\
(b) Seefläche & 1\E & $3.2 \cdot 625\,000\,000\,\text{cm}^2 = 2\,000\,000\,000\,\text{cm}^2 = 200\,000\,\text{m}^2 = 0.2\,\text{km}^2$\\
(c) neue Karte & 1 & Längen auf der Karte halb so lang ($\tfrac{25\,000}{50\,000} = \tfrac{1}{2}$), Flächen ein Viertel: $3.2 : 4 = 0.8\,\text{cm}^2$ (gleichwertig: $0.2\,\text{km}^2 : 50\,000^2$)\\
\end{raster}
\fehler{Fläche mit $25\,000$ statt $25\,000^2$ ($80\,000\,\text{cm}^2 = 8\,\text{m}^2$): beide (b)-Zeilen 0. Rechnet (c) mit derselben Verwechslung ($3.2 : 2 = 1.6\,\text{cm}^2$), ist das derselbe Fehler noch einmal, kein Folgefehler: (c) 0 $\to$ 1 von 4\,P.
Umrechnung $\text{cm}^2 \to \text{km}^2$ falsch (z.\,B. $20\,\text{km}^2$ oder $0.002\,\text{km}^2$) oder nicht umgerechnet ($2\,000\,000\,000\,\text{cm}^2$, $200\,000\,\text{m}^2$): (b)-Seefläche 0, Faktorzeile zählt. Ebenso (a): $460\,000\,\text{cm}$ oder $4600\,\text{m}$ statt km gibt (a) 0.}
\end{lpaufgabe}

\begin{lpaufgabe}{G5}{Blumenbeet}{3 P}
\begin{raster}
(a) Faktor & 1 & Fläche mal $3$, also $k^2 = 3$, $k = \sqrt{3} \approx 1.73$\\
(a) Durchmesser & 1\E & $1.6 \cdot \sqrt{3} \approx 2.77\,\text{m}$ (mit $1.73$ gerechnet: $2.77$)\\
(b) Umfang & 1 & Der Umfang ist eine Länge: Faktor $k = \sqrt{3} \approx 1.73$. Folgepunkt: das eigene $k$, wenn es als Längenfaktor genommen wird.\\
\end{raster}
\fehler{$k = 3$ (Flächenfaktor als Längenfaktor), Durchmesser $4.8\,\text{m}$: beide (a)-Zeilen 0; (b) mit dem Faktor $3$ zählt als Folgepunkt — der Schritt in (b), «der Umfang wächst mit dem Längenfaktor», ist richtig, nur der Faktor kommt aus (a) $\to$ 1 von 3\,P.
(b) «mal $3$», obwohl in (a) $k = \sqrt{3}$ steht: (b) 0.}
\end{lpaufgabe}

\begin{lpaufgabe}{G6}{Zwei Dreiecke}{4 P}
\begin{raster}
(a) ähnlich & 1 & Dritter Winkel bei $DEF$: $\angle F = 180^\circ - 41^\circ - 76^\circ = 63^\circ$; bei $XYZ$: $\angle Y = 180^\circ - 63^\circ - 41^\circ = 76^\circ$. Beide haben $41^\circ$, $63^\circ$, $76^\circ$: ähnlich (WW).\\
(a) Ecken & 1\A & $X \leftrightarrow F$ ($63^\circ$), $Y \leftrightarrow E$ ($76^\circ$), $Z \leftrightarrow D$ ($41^\circ$). Alle drei nötig.\\
(b) Laras Fehler & 1\A & Sie ordnet nach den Buchstaben zu statt nach den Winkeln. $YZ$ liegt dem Winkel $63^\circ$ (bei $X$) gegenüber und entspricht $DE$ (gegenüber $F$), nicht $EF$.\\
(b) $\overline{XY}$ & 1\E & $k = \tfrac{\overline{YZ}}{\overline{DE}} = \tfrac{6}{8} = 0.75$; $XY$ liegt $41^\circ$ gegenüber und entspricht $EF$: $\overline{XY} = 0.75 \cdot 5.89 \approx 4.42\,\text{cm}$ (aus den ungerundeten Seiten ebenfalls $4.42$)\\
\end{raster}
\fehler{Laras Rechnung übernommen ($\overline{XY} \approx 8.15$): Ecken-, Fehler- und $\overline{XY}$-Zeile 0 $\to$ 1 von 4\,P.
$k = \tfrac{8}{6}$ (verkehrt), $\overline{XY} = 5.89 \cdot \tfrac{8}{6} \approx 7.85$: $\overline{XY}$-Zeile 0.
Richtiges $k$, aber $XY$ zu $FD$ zugeordnet ($0.75 \cdot 8.71 \approx 6.53$): $\overline{XY}$-Zeile 0 — derselbe Zuordnungsfehler, kein Folgefehler.
«Lara hat falsch gerechnet» ohne Grund: Fehlerzeile 0.}
\end{lpaufgabe}

\begin{lpaufgabe}{G7}{Rechtwinkliges Dreieck}{3 P}
\begin{raster}
(a) $p$ & 1\E & Kathetensatz $a^2 = p \cdot c$: $p = \tfrac{7.5^2}{12.5} = \tfrac{56.25}{12.5} = 4.5\,\text{cm}$\\
(a) $q$ & 1 & $q = c - p = 8\,\text{cm}$. Folgepunkt: $c$ minus das eigene $p$.\\
(b) $h$ & 1 & Höhensatz $h^2 = p \cdot q = 36$, $h = 6\,\text{cm}$ (gleichwertig: Pythagoras im Teildreieck, $h = \sqrt{7.5^2 - 4.5^2}$). Folgepunkt: richtig mit den eigenen $p$, $q$.\\
\end{raster}
\fehler{$p = \tfrac{7.5}{12.5} = 0.6$ (Quadrat vergessen): $p$-Zeile 0; $q = 11.9$ und $h = \sqrt{0.6 \cdot 11.9} \approx 2.67$ zählen als Folgepunkte $\to$ 2 von 3\,P.
$p$ und $q$ vertauscht ($p = 8$, $q = 4.5$): $p$-Zeile 0, die übrigen zählen ($h = 6$ bleibt).}
\end{lpaufgabe}

\needspace{14\baselineskip}
\section*{4 · Selbsteinschätzung}
\begin{tabularx}{\linewidth}{@{}l X@{}}\toprule
22 – 25 P & Die geprüften Teile sitzen. Wo du Punkte verloren hast: das Kapitel dieser Aufgabe nochmals (Zuordnung unten).\\
17 – 21 P & Das schwächste Kapitel nochmals: Tüfteln und Übungen des Kapitels, in dem du die meisten Punkte verloren hast.\\
11 – 16 P & Zurück zu den Kapiteln aller Aufgaben, in denen du Punkte verloren hast.\\
0 – 10 P & Zurück zu Kapitel 1 und von dort der Reihe nach weiter.\\\midrule
\multicolumn{2}{@{}p{\linewidth}@{}}{\small\textbf{Aufgabe $\to$ Kapitel:} G1 $\to$ Kapitel 1 (Zentrische Streckung); G2, G3 $\to$ Kapitel 2 (Strahlensätze); G4, G5 $\to$ Kapitel 3 (Ähnliche Figuren); G6, G7 $\to$ Kapitel 4 (Ähnliche Dreiecke)}\\\bottomrule
\end{tabularx}
\end{document}
